{\color{ink}\bfseries 例 2-17}\par\vspace{1bp}Need按P0至P4排列为(0,\allowbreak{}0,\allowbreak{}1)、(0,\allowbreak{}7,\allowbreak{}5)、(1,\allowbreak{}0,\allowbreak{}0)、(0,\allowbreak{}6,\allowbreak{}2)、(0,\allowbreak{}6,\allowbreak{}4)。\par I可取P2\(\to \)P0\(\to \)P1\(\to \)P3\(\to \)P4，Work依次为(1,\allowbreak{}4,\allowbreak{}0)\(\to \)(2,\allowbreak{}7,\allowbreak{}5)\(\to \)(2,\allowbreak{}7,\allowbreak{}8)\(\to \)(3,\allowbreak{}7,\allowbreak{}8)\(\to \)(3,\allowbreak{}7,\allowbreak{}10)\(\to \)(3,\allowbreak{}7,\allowbreak{}11)。\par II完成P0、P3、P4后只能到(0,\allowbreak{}6,\allowbreak{}8)：P1缺B，P2缺A；III完成P0、P2后到(2,\allowbreak{}4,\allowbreak{}9)，其余进程均至少还需6个B；IV完成P0后为(0,\allowbreak{}4,\allowbreak{}10)，其余进程仍不满足。故只有I安全，选C；能启动某个进程只证明存在可行前缀。
